% GENERATED FILE. Edit canonical lessons, then run npm run generate:books.
% canonical-source-hash: fnv1a64:8b74a7b71bd99565
% canonical-lessons: ML-C309-kuraykkuka2, ML-C309-tiriccariyuka, ML-C309-puttuka, ML-C309-vivarikkuka, ML-C309-olivakkuka

\chapter{To reduce, To recognize, To lock, To describe, To avoid}
\label{ch:ml-to-reduce-to-recognize-to-lock-to-describe-to-avoid}
\hlchaptermodality{309}

\begin{chapteropening}
\textbf{By the end of this chapter:} \emph{I can say to reduce, to recognize, to lock, to describe and to avoid.}

Complete the last lesson of chapter 309: I can say to reduce, to recognize, to lock, to describe and to avoid.
\end{chapteropening}

\section[\texorpdfstring{\ml{കുറയ്ക്കുക} (kuṟaykkuka)}{കുറയ്ക്കുക (kuṟaykkuka)}]{\ml{കുറയ്ക്കുക} (kuṟaykkuka) — to reduce}

\label{lesson:ML-C309-kuraykkuka2}



\begin{quote}
Before the new one: say the Malayalam for to take a photo, then the Malayalam for to type.
\end{quote}

\subsection*{You'll want to know: \ml{കുറയ്ക്കുക}}
\textbf{\ml{കുറയ്ക്കുക}} — \emph{kuṟaykkuka} — \textquotedblleft{}to reduce\textquotedblright{}.

\begin{cousinweb}
\textbf{\ml{കുറയ്ക്കുക}} (\emph{kuṟaykkuka}) is native Dravidian, the causative of the verb \textquotedblleft{}to dwindle\textquotedblright{}: \textquotedblleft{}make less\textquotedblright{}; DEDR 1851 lists the root with Tamil \emph{kuṟu}.
\end{cousinweb}

\begin{grammarlens}[title={What you've built}]
One more word for this chapter.
\end{grammarlens}

\subsection*{Your turn}
\begin{itemize}
\raggedright
  \item \emph{Say it:} \emph{kuṟaykkuka}
  \item \emph{Say it:} \emph{kuṟaykkuka}, once more
  \item \emph{Say it:} say \emph{kuṟaykkuka}
  \item \emph{Recall:} say the Malayalam for to take a photo, then the Malayalam for to type, then say \emph{kuṟaykkuka} again
\end{itemize}

\begin{tcolorbox}[breakable,colback=teal!4,colframe=teal!35!black,title={Before you move on}]
What does \ml{കുറയ്ക്കുക} mean? (\textquotedblleft{}To reduce\textquotedblright{}.) Say it once more.
\end{tcolorbox}
\section[\texorpdfstring{\ml{തിരിച്ചറിയുക} (tiriccaṟiyuka)}{തിരിച്ചറിയുക (tiriccaṟiyuka)}]{\ml{തിരിച്ചറിയുക} (tiriccaṟiyuka) — to recognize}

\label{lesson:ML-C309-tiriccariyuka}



\begin{quote}
Before the new one: say the Malayalam for to type, then the Malayalam for to reduce.
\end{quote}

\subsection*{You'll want to know: \ml{തിരിച്ചറിയുക}}
\textbf{\ml{തിരിച്ചറിയുക}} — \emph{tiriccaṟiyuka} — \textquotedblleft{}to recognize\textquotedblright{}.

\begin{cousinweb}
\textbf{\ml{തിരിച്ചറിയുക}} (\emph{tiriccaṟiyuka}) is native Dravidian, \textquotedblleft{}to know again\textquotedblright{}: \emph{tiriccu}, \textquotedblleft{}turning back\textquotedblright{} (DEDR 3246, with Tamil \emph{tiri}), and \emph{aṟiyuka}, \textquotedblleft{}to know\textquotedblright{} (DEDR 314).
\end{cousinweb}

\begin{grammarlens}[title={What you've built}]
One more word for this chapter.
\end{grammarlens}

\subsection*{Your turn}
\begin{itemize}
\raggedright
  \item \emph{Say it:} \emph{tiriccaṟiyuka}
  \item \emph{Say it:} \emph{tiriccaṟiyuka}, once more
  \item \emph{Say it:} say \emph{tiriccaṟiyuka}
  \item \emph{Recall:} say the Malayalam for to type, then the Malayalam for to reduce, then say \emph{tiriccaṟiyuka} again
\end{itemize}

\begin{tcolorbox}[breakable,colback=teal!4,colframe=teal!35!black,title={Before you move on}]
What does \ml{തിരിച്ചറിയുക} mean? (\textquotedblleft{}To recognize\textquotedblright{}.) Say it once more.
\end{tcolorbox}
\section[\texorpdfstring{\ml{പൂട്ടുക} (pūṭṭuka)}{പൂട്ടുക (pūṭṭuka)}]{\ml{പൂട്ടുക} (pūṭṭuka) — to lock}

\label{lesson:ML-C309-puttuka}



\begin{quote}
Before the new one: say the Malayalam for to reduce, then the Malayalam for to recognize.
\end{quote}

\subsection*{You'll want to know: \ml{പൂട്ടുക}}
\textbf{\ml{പൂട്ടുക}} — \emph{pūṭṭuka} — \textquotedblleft{}to lock\textquotedblright{}.

\begin{cousinweb}
\textbf{\ml{പൂട്ടുക}} (\emph{pūṭṭuka}) is native Dravidian, the causative of the verb \textquotedblleft{}to be closed, be yoked\textquotedblright{}; DEDR 4361 lists the root with Tamil \emph{pūṇ}, Kannada \emph{pūṇ} and Telugu \emph{pūnu}.
\end{cousinweb}

\begin{grammarlens}[title={What you've built}]
One more word for this chapter.
\end{grammarlens}

\subsection*{Your turn}
\begin{itemize}
\raggedright
  \item \emph{Say it:} \emph{pūṭṭuka}
  \item \emph{Say it:} \emph{pūṭṭuka}, once more
  \item \emph{Say it:} say \emph{pūṭṭuka}
  \item \emph{Recall:} say the Malayalam for to reduce, then the Malayalam for to recognize, then say \emph{pūṭṭuka} again
\end{itemize}

\begin{tcolorbox}[breakable,colback=teal!4,colframe=teal!35!black,title={Before you move on}]
What does \ml{പൂട്ടുക} mean? (\textquotedblleft{}To lock\textquotedblright{}.) Say it once more.
\end{tcolorbox}
\section[\texorpdfstring{\ml{വിവരിക്കുക} (vivarikkuka)}{വിവരിക്കുക (vivarikkuka)}]{\ml{വിവരിക്കുക} (vivarikkuka) — to describe, to explain}

\label{lesson:ML-C309-vivarikkuka}



\begin{quote}
Before the new one: say the Malayalam for to recognize, then the Malayalam for to lock.
\end{quote}

\subsection*{You'll want to know: \ml{വിവരിക്കുക}}
\textbf{\ml{വിവരിക്കുക}} — \emph{vivarikkuka} — \textquotedblleft{}to describe, to explain\textquotedblright{}.

\begin{cousinweb}
\textbf{\ml{വിവരിക്കുക}} (\emph{vivarikkuka}) is built on Sanskrit \emph{vivaraṇa}, \textquotedblleft{}opening up, explanation\textquotedblright{}, from \emph{vivara}, \textquotedblleft{}an opening\textquotedblright{} (CDIAL 11915), with the ending that makes a Malayalam verb.
\end{cousinweb}

\begin{grammarlens}[title={What you've built}]
One more word for this chapter.
\end{grammarlens}

\subsection*{Your turn}
\begin{itemize}
\raggedright
  \item \emph{Say it:} \emph{vivarikkuka}
  \item \emph{Say it:} \emph{vivarikkuka}, once more
  \item \emph{Say it:} say \emph{vivarikkuka}
  \item \emph{Recall:} say the Malayalam for to recognize, then the Malayalam for to lock, then say \emph{vivarikkuka} again
\end{itemize}

\begin{tcolorbox}[breakable,colback=teal!4,colframe=teal!35!black,title={Before you move on}]
What does \ml{വിവരിക്കുക} mean? (\textquotedblleft{}To describe, to explain\textquotedblright{}.) Say it once more.
\end{tcolorbox}
\section[\texorpdfstring{\ml{ഒഴിവാക്കുക} (oḻivākkuka)}{ഒഴിവാക്കുക (oḻivākkuka)}]{\ml{ഒഴിവാക്കുക} (oḻivākkuka) — to avoid}

\label{lesson:ML-C309-olivakkuka}



\begin{quote}
Before the new one: say the Malayalam for to lock, then the Malayalam for to describe.
\end{quote}

\subsection*{You'll want to know: \ml{ഒഴിവാക്കുക}}
\textbf{\ml{ഒഴിവാക്കുക}} — \emph{oḻivākkuka} — \textquotedblleft{}to avoid\textquotedblright{}.

\begin{cousinweb}
\textbf{\ml{ഒഴിവാക്കുക}} (\emph{oḻivākkuka}) is native Dravidian, \textquotedblleft{}make it a vacancy\textquotedblright{}: \emph{oḻivu}, \textquotedblleft{}being free, vacancy\textquotedblright{} (DEDR 1009), and \emph{ākkuka}, \textquotedblleft{}to make\textquotedblright{} (DEDR 333). DEDR 333 compares Sanskrit \emph{ām}.
\end{cousinweb}

\begin{grammarlens}[title={What you've built}]
One more word for this chapter.
\end{grammarlens}

\subsection*{Your turn}
\begin{itemize}
\raggedright
  \item \emph{Say it:} \emph{oḻivākkuka}
  \item \emph{Say it:} \emph{oḻivākkuka}, once more
  \item \emph{Say it:} say \emph{oḻivākkuka}
  \item \emph{Recall:} say the Malayalam for to lock, then the Malayalam for to describe, then say \emph{oḻivākkuka} again
\end{itemize}

\begin{tcolorbox}[breakable,colback=teal!4,colframe=teal!35!black,title={Before you move on}]
What does \ml{ഒഴിവാക്കുക} mean? (\textquotedblleft{}To avoid\textquotedblright{}.) Say it once more.
\end{tcolorbox}
